% p2-11-limitations

\section{P2 §11. Limitations of the Present Work}

11. Limitations of the Present Work

    Figure (fig-p2-ch11-limitations). Limitations triptych --- Toy-model caveat / No derivation / Data
                                                   gaps.

  The following limitations are acknowledged explicitly and without mitigation.

  No new derivations. This paper derives no new field-theory results. The E8/Toda
  material is reviewed from the literature. The sRG model and operator framework are
  formal constructions without derived content.

  The mechanism connection is asserted, not derived. The link between the E8/Toda $\varphi
  occurrences and the Paper 1 compressibility data is a research hypothesis, not a
  demonstrated fact.

  Selectivity risk. The focus on $\varphi$ rather than other algebraic numbers of similar degree
  (2, 3, 6) reflects the specific output of Paper 1 but also introduces the possibility of
  post-hoc selection. WP5 is designed to address this, but until it is completed, selectivity
  risk remains. The flos34 catalogue was developed with $\varphi$ as a preferred basis; using it
  in WP5 without blinding introduces circular risk that must be managed by the
  experimental design.

  2D vs. 4D gap. All established $\varphi$ field-theory results (Zamolodchikov E8, Toda/H3, H4)
  are in 2D. The Standard Model is a 4D theory. The gap between 2D integrable models
  and 4D non-integrable gauge theory is not bridged by anything in this paper. This gap is
  the single largest obstacle to the programme.

  The $\pi^{2}$-incommensurability obstacle. As shown in Section 5.3, the one-loop beta
  function of phi4 theory introduces pi2, which lies outside Q(5) by algebraic
  independence (Nesterenko's theorem). This structural fact means that $\varphi$-closure of RG
  trajectories --- if it exists --- cannot be a property of individual coupling beta functions at
  standard perturbative order; it must emerge in ratios or at fixed points.

  DSI mechanism is entirely conjectural. The log-periodic correction formula (21) is a
  template with no derived amplitude. No prediction of a specific numerical value is
  made. The two $\alpha$-$\varphi$ candidates --- alpha$\varphi$,log = ln($\varphi$2)/pi approx 0.306 and alpha$\varphi$,alg =
  $\varphi$-3/2 approx 0.118 --- are not used in the DSI template because they have not been
  derived from first principles; their disambiguation is a prerequisite for Paper 3. These
  two quantities must not be conflated.

  Coq coverage is partial. The Coq proof infrastructure covers identity (1), a
  representative catalogue layer, and several invariants in the IGLA training bundle, but
  does not cover the $\varphi$-structured operator hypothesis or the symbolic RG model. A local
  source audit of the available t27 checkout records the topic-specific t27/proofs/trinity/
  layer as 13 Coq files, 281 theorem-like declarations, 122 source-level Qed. markers, 32
  source-level Admitted. markers, and 0 Abort. markers. The broader t27/proofs/ tree
  records 18 Coq proof files, 296 theorem-like declarations, 127 Qed. markers, 32
  Admitted. markers, and 0 Abort. markers; the older t27/coq/ kernel records 10 Coq
  files, 73 declarations, 35 Qed. markers, 0 Admitted. markers, and 11 Abort. markers.
  The catalogue aggregation file t27/proofs/trinity/Catalog42.v contains 13 source-level
  Qed. markers and 0 Admitted. markers, and defines a 28-obligation representative
  verification aggregate. Therefore the 42-entry catalogue should be cited as a 42-row
  formula catalogue plus a larger Coq proof layer, not as forty-two independent row-wise
  Coq theorems. None of those obligations affect the physical-mechanism claims of the
  present paper, which are conjectural regardless.

  No data. This paper contains no new experimental data. It is entirely theoretical and
  programmatic.

\subsection{References}

- Coldea, R. et al. "Quantum Criticality in an Ising Chain: Experimental Evidence for Emergent E8 Symmetry." \textit{Science} 327 (2010) 177. DOI \href{https://doi.org/10.1126/science.1180085}{10.1126/science.1180085}.
- Wu, J. et al. "E8 Symmetry in CoN$b_{2}$$O_{6}$ at Higher Resolution." \textit{Phys. Rev. B} 108, L060403 (2023). DOI \href{https://doi.org/10.1103/PhysRevB.108.L060403}{10.1103/PhysRevB.108.L060403}.
- Fring, A.; Korff, C. "Affine Toda field theories related to Coxeter groups of non-crystallographic type." \textit{Nucl. Phys. B} 729 (2005) 361. \href{https://arxiv.org/abs/hep-th/0506226}{arXiv:hep-th/0506226}. DOI \href{https://doi.org/10.1016/j.nuclphysb.2005.08.044}{10.1016/j.nuclphysb.2005.08.044}.
- Pellis, S. "The Fine-Structure Constant and the Golden Ratio." \href{https://vixra.org/pdf/2110.0117v4.pdf}{viXra 2110.0117 v4}.
